\documentclass[12pt,reqno]{article}

\usepackage[usenames]{color}
\usepackage{amssymb}
\usepackage{amsmath}
\usepackage{amsthm}
\usepackage{amsfonts}
\usepackage{amscd}
\usepackage{graphicx}
\usepackage{xcolor}
\usepackage{float}
\usepackage{booktabs}


\usepackage[colorlinks=true,
linkcolor=webgreen,
filecolor=webbrown,
citecolor=webgreen]{hyperref}

\definecolor{webgreen}{rgb}{0,.5,0}
\definecolor{webbrown}{rgb}{.6,0,0}


\usepackage{fullpage}

%\usepackage{psfig}
\usepackage{graphics}
\usepackage{latexsym}
\usepackage{epsf}
\usepackage{breakurl}

\setlength{\textwidth}{6.5in}
\setlength{\oddsidemargin}{.1in}
\setlength{\evensidemargin}{.1in}
\setlength{\topmargin}{-.1in}
\setlength{\textheight}{8.4in}

\newcommand{\seqnum}[1]{\href{https://oeis.org/#1}{\rm \underline{#1}}}

\newcommand{\Z}{\mathbb Z}
\newcommand{\N}{\mathbb N}
\newcommand{\word}[1]{\mathtt{#1}}
\newcommand{\vp}{\nu_3}

\newcommand{\Q}{\mathbb{Q}}
\newcommand{\wt}{\mathbf{t}}
\newcommand{\ws}{\mathbf{s}}
\newcommand{\ww}{\mathbf{w}}

\DeclareMathOperator{\Par}{\psi}

\begin{document}


\theoremstyle{plain}
\newtheorem{theorem}{Theorem}
\newtheorem{corollary}[theorem]{Corollary}
\newtheorem{lemma}[theorem]{Lemma}
\newtheorem{proposition}[theorem]{Proposition}

\theoremstyle{definition}
\newtheorem{definition}[theorem]{Definition}
\newtheorem{example}[theorem]{Example}
\newtheorem{conjecture}[theorem]{Conjecture}

\theoremstyle{remark}
\newtheorem{remark}[theorem]{Remark}

\author{S. Cambie, E. Kalviainen, J. Shallit}

\title{Gerver-Ramsey Theorems in Small Dimensions}

\maketitle

\begin{abstract}

\end{abstract}

\section{Introduction}

history of the problem

restriction to unit vectors

results so far (Lidbetter)

summary of our results

notational conventions:  Throughout, $\N=\{0,1,2,\ldots\}$, and word and vertex indices start
at zero.

connection to weak abelian powers:  Avoiding $k+1$ collinear vertices is equivalent to avoiding weak abelian
$k$th powers: if $0\leq n_0<\cdots<n_k$ are collinear vertices, their
consecutive displacements are positive multiples of a common vector;
division by their coordinate sums gives the same normalized Parikh
vector for all $k$ intervening nonempty blocks. Conversely, such $k$
blocks give $k+1$ collinear boundary vertices.

It's possible that in $\N^4$ there are infinite walks, with unit vector steps, with no three collinear points.  We have constructed a finite walk of length $2600$ with this property.

It's possible that in $\N^3$ there are infinite walks, with unit vector steps, with no four collinear points.  We have constructed a finite walk of length $32768$ with this property.


\section{Dimensions 6 and up}

The new Cambie-Kalviainen result shows you can avoid three collinear points in dimension 3 with steps of 6 different vectors.  From this it easily follows that in dimension 6 one can avoid 3 collinear points using only the unit vectors in $\N^6$ as steps, using one unit vector to represent each of the 6 Cambie-Kalviainen  vectors.



\section{Dimension 4}
\label{dim4}

For a finite word
$v$ over $\Sigma=\{0,1,2,3\}$, let $|v|$ denote its length, let $|v|_j$
denote the number of occurrences of the letter $j$, and define its Parikh
vector by
\[
  \Par(v)=(|v|_0,|v|_1,|v|_2,|v|_3)^{\mathsf T}\in\Z_{\geq0}^{4}.
\]
A \emph{weak abelian cube} is a concatenation $xyz$ of three nonempty
finite words such that
\[
  \frac{\Par(x)}{|x|}
  =\frac{\Par(y)}{|y|}
  =\frac{\Par(z)}{|z|}.
\]
The three block lengths need not be equal. A factor of an infinite word
means a finite contiguous block.

The construction below is a ternary adaptation of the valuation method
used by Cambie and Kalviainen~\cite{CK} for a finite-step lattice walk
with no three collinear vertices.

\begin{theorem}\label{thm:main}
Let $\mathcal A=\{A,B,C,D,E\}$. Define a substitution
$\tau:\mathcal A^*\to\mathcal A^*$ and a letter-to-letter coding
$h:\mathcal A^*\to\Sigma^*$ by
\[
\begin{array}{c|ccccc}
a&A&B&C&D&E\\ \hline
\tau(a)&AB&AACA&ADE&AACCE&ADCCA\\
h(a)&0&1&2&3&3
\end{array}
\]
and extend both maps to words by concatenation. Then
\[
  \ww=h\bigl(\tau^\infty(A)\bigr)
\]
is an infinite word over $\Sigma$ containing no weak abelian cube.
\end{theorem}

The limit $\tau^\infty(A)$ is well defined: $\tau(A)$ starts with $A$,
so $\tau^k(A)$ is a prefix of $\tau^{k+1}(A)$, and all images have
length at least $2$, so these prefix lengths tend to infinity. Since $h$
maps every auxiliary letter to one letter, its image is also infinite.
For reference, the word begins
\[
  \ww=\word{0100200101033010100200100200100223}\cdots.
\]
All four output letters occur: $A,B,C$ occur in $\tau^2(A)=ABAACA$,
and $D,E$ occur in $\tau^3(A)$.

We first identify the blocks encoded by the auxiliary letters. Define
another substitution, this time on $\Sigma$, by
\begin{equation}\label{eq:sigma}
  \sigma(0)=\word{010},\qquad
  \sigma(1)=\word{232},\qquad
  \sigma(2)=\word{101},\qquad
  \sigma(3)=\word{323},
\end{equation}
and put $\wt=\sigma^\infty(0)$. This fixed point exists because
$\sigma(0)$ starts with $0$ and every image has length $3$.
Define $R:\mathcal A^*\to\Sigma^*$ by the following nonempty images:
\begin{equation}\label{eq:R}
\begin{array}{c|ccccc}
a&A&B&C&D&E\\ \hline
R(a)&\word{01}&\word{0232}&\word{013231}&
\word{02323231}&\word{01323232}.
\end{array}
\end{equation}

\begin{lemma}\label{lem:returns}
Writing $\ws=s_0s_1s_2\cdots=\tau^\infty(A)$, we have
\[
  R(\ws)=R(s_0)R(s_1)R(s_2)\cdots=\wt.
\]
Moreover, the starts of these blocks are exactly the positions at which
$\wt$ has the letter $0$.
\end{lemma}

\begin{proof}
The five required substitution identities are as follows. Spaces on the
right separate the successive $R$-images.
\begin{align*}
\sigma(R(A))&=\word{01\ 0232}
             =R(A)R(B),\\
\sigma(R(B))&=\word{01\ 01\ 013231\ 01}
             =R(A)R(A)R(C)R(A),\\
\sigma(R(C))&=\word{01\ 02323231\ 01323232}
             =R(A)R(D)R(E),\\
\sigma(R(D))&=\word{01\ 01\ 013231\ 013231\ 01323232}
             =R(A)R(A)R(C)R(C)R(E),\\
\sigma(R(E))&=\word{01\ 02323231\ 013231\ 013231\ 01}
             =R(A)R(D)R(C)R(C)R(A).
\end{align*}
Thus $\sigma\circ R=R\circ\tau$ on every letter, and hence on every
finite word. Induction on $k$ gives
\[
  R(\tau^k(A))=\sigma^k(R(A))=\sigma^k(\word{01}).
\]
The word on the right starts with $\sigma^k(0)$, whose length is $3^k$.
Consequently, for every fixed prefix length, these prefixes eventually
agree with $\wt$. On the left, the nested words converge to $R(\ws)$.
This proves $R(\ws)=\wt$.

Each of the five words in~\eqref{eq:R} begins with $0$ and has no other
occurrence of $0$. The concatenation is infinite and each block is
nonempty. Its block starts are therefore precisely all the zero
positions of $\wt$.
\end{proof}

We next associate a lattice walk with $\wt=t_0t_1t_2\cdots$. For each
$n\in\N$, define $(a_n,b_n)$ from $t_n$ by
\begin{equation}\label{eq:states}
\begin{array}{c|rrrr}
t_n&0&1&2&3\\ \hline
a_n&1&1&-1&-1\\
b_n&1&-1&1&-1
\end{array}
\end{equation}
and define integer prefix sums and vertices by
\begin{equation}\label{eq:vertices}
  X_n=\sum_{j=0}^{n-1}a_j,\qquad
  Y_n=\sum_{j=0}^{n-1}b_j,\qquad
  P_n=(X_n,Y_n,n)\in\Z^3.
\end{equation}
An empty sum is zero, so $P_0=(0,0,0)$. Finally, put
\[
  F(U,V)=U^2-3V^2.
\]
For a nonzero integer $q$, $\nu_3(q)$ denotes the largest nonnegative
integer $e$ such that $3^e$ divides $q$; the sign of $q$ is irrelevant.
Extend this to nonzero rational numbers by
$\nu_3(p/q)=\nu_3(p)-\nu_3(q)$. In particular,
$\nu_3(uv)=\nu_3(u)+\nu_3(v)$ for nonzero rational $u,v$.

\begin{lemma}\label{lem:recursion}
For every $n\in\N$ and $r\in\{0,1,2\}$,
\begin{align}
  a_{3n+r}&=b_n,& b_{3n+r}&=(-1)^r a_n,\label{eq:state-recursion}\\
  X_{3n+r}&=3Y_n+r b_n,&
  Y_{3n+r}&=X_n+\varepsilon_r a_n,\label{eq:sum-recursion}
\end{align}
where $\varepsilon_0=0$, $\varepsilon_1=1$, and $\varepsilon_2=0$.
\end{lemma}

\begin{proof}
For an input state $(a_n,b_n)$, the three output states in its
$\sigma$-image, as read from~\eqref{eq:sigma} and~\eqref{eq:states}, are
\[
  (b_n,a_n),\qquad(b_n,-a_n),\qquad(b_n,a_n).
\]
This proves~\eqref{eq:state-recursion}. The sum of their first
coordinates is $3b_n$, and the sum of their second coordinates is
$a_n$. Thus $X_{3n}=3Y_n$ and $Y_{3n}=X_n$. Within the next block, the
first $r$ first coordinates sum to $r b_n$, while the first $r$ second
coordinates sum to $0,a_n,0$ for $r=0,1,2$, respectively. Adding these
partial sums proves~\eqref{eq:sum-recursion}.
\end{proof}

\begin{lemma}[Equal-state valuation identity]\label{lem:valuation}
If $0\leq m<n$ and $t_m=t_n$, then
\[
  F(X_n-X_m,Y_n-Y_m)\ne0
\]
and
\begin{equation}\label{eq:valuation}
  \nu_3\bigl(F(X_n-X_m,Y_n-Y_m)\bigr)=\nu_3(n-m).
\end{equation}
\end{lemma}

\begin{proof}
Write $m=3u+r$ and $n=3v+s$, with $r,s\in\{0,1,2\}$. We separate
the cases $r\ne s$ and $r=s$.

Suppose first that $r\ne s$. The endpoint states agree, so
$a_m=a_n$. Equation~\eqref{eq:state-recursion} implies
\[
  b_u=a_m=a_n=b_v.
\]
Consequently, by~\eqref{eq:sum-recursion},
\[
  X_n-X_m
  =3(Y_v-Y_u)+s b_v-r b_u
  \equiv(s-r)a_m\pmod3.
\]
Here $a_m\in\{-1,1\}$, and $s-r$ is not divisible by $3$.
Therefore $X_n-X_m$ is not divisible by $3$, and
\[
  F(X_n-X_m,Y_n-Y_m)\equiv(X_n-X_m)^2\not\equiv0\pmod3.
\]
The quadratic expression is nonzero and has valuation zero. Also
$3\nmid n-m$, so~\eqref{eq:valuation} holds in this case.

Now suppose $r=s$. Equality of the first coordinates of the endpoint
states gives $b_u=b_v$, and equality of the second coordinates gives
$(-1)^r a_u=(-1)^r a_v$, hence $a_u=a_v$. The correspondence in
\eqref{eq:states} is one-to-one, so $t_u=t_v$. In particular, the
correction terms in~\eqref{eq:sum-recursion} cancel, giving
\[
  X_n-X_m=3(Y_v-Y_u),\qquad
  Y_n-Y_m=X_v-X_u.
\]
It follows that
\begin{align}
 F(X_n-X_m,Y_n-Y_m)
 &=9(Y_v-Y_u)^2-3(X_v-X_u)^2\notag\\
 &=-3F(X_v-X_u,Y_v-Y_u).\label{eq:descent}
\end{align}
At the same time, $n-m=3(v-u)$, and $u<v$.

Let $e=\nu_3(n-m)$. Apply this second case exactly $e$ times. At each
step the endpoint states still agree, the indices remain distinct and
ordered, and their difference is divided by $3$. We arrive at indices
$m'<n'$ such that $t_{m'}=t_{n'}$ and $3\nmid n'-m'$, with
\[
  F(X_n-X_m,Y_n-Y_m)
  =(-3)^e F(X_{n'}-X_{m'},Y_{n'}-Y_{m'}).
\]
The first case shows that the final factor is nonzero and has valuation
zero. The original expression is therefore nonzero and has valuation
$e$, as required.
\end{proof}

\begin{proposition}\label{prop:no-four}
No four distinct vertices in the set
\[
  \{P_n:n\in\N,\ t_n=0\}
\]
are collinear.
\end{proposition}

\begin{proof}
Suppose that $P_{n_0},P_{n_1},P_{n_2},P_{n_3}$ are collinear, where
$n_0<n_1<n_2<n_3$ and $t_{n_i}=0$ for each $i$. Their third
coordinates are distinct, so the line has well-defined rational slopes
\[
  \alpha=\frac{X_{n_1}-X_{n_0}}{n_1-n_0},\qquad
  \beta=\frac{Y_{n_1}-Y_{n_0}}{n_1-n_0}.
\]
For every $0\leq i<j\leq3$, collinearity gives
\[
 X_{n_j}-X_{n_i}=\alpha(n_j-n_i),\qquad
 Y_{n_j}-Y_{n_i}=\beta(n_j-n_i).
\]
Put $c=\alpha^2-3\beta^2$. Applying
Lemma~\ref{lem:valuation} to the pair $(n_0,n_1)$ shows that $c\ne0$.
For every pair $(i,j)$, the same lemma now gives
\begin{align*}
 \nu_3(n_j-n_i)
 &=\nu_3\bigl(c(n_j-n_i)^2\bigr)\\
 &=\nu_3(c)+2\nu_3(n_j-n_i).
\end{align*}
Hence all six positive differences $n_j-n_i$ have the same valuation
\[
  e=-\nu_3(c)\in\Z_{\geq0}.
\]
The four numbers
\[
  q_i=\frac{n_i-n_0}{3^e}\qquad(0\leq i\leq3)
\]
are integers. For $i<j$, the difference
$q_j-q_i=(n_j-n_i)/3^e$ is not divisible by $3$. Thus the four $q_i$
would be pairwise distinct modulo $3$. There are only three residue
classes modulo $3$, a contradiction.
\end{proof}

We can now relate the geometry to the normalized Parikh vectors of
$\ww$. For a finite word $v$ over $\Sigma$, define its displacement by
\[
  \Delta(v)=\bigl(|v|_0+|v|_1-|v|_2-|v|_3,\,
  |v|_0-|v|_1+|v|_2-|v|_3,\,|v|\bigr)^{\mathsf T}.
\]
By~\eqref{eq:states} and~\eqref{eq:vertices}, if $v$ is the factor
$t_m t_{m+1}\cdots t_{n-1}$, then $\Delta(v)=P_n-P_m$.
The return blocks have the following Parikh vectors and displacements:
\[
\begin{array}{cclc}
\toprule
a & h(a) & \Par(R(a))^{\mathsf T} & \Delta(R(a))^{\mathsf T}\\
\midrule
A&0&(1,1,0,0)&(2,0,2)\\
B&1&(1,0,2,1)&(-2,2,4)\\
C&2&(1,2,1,2)&(0,-2,6)\\
D&3&(1,1,3,3)&(-4,0,8)\\
E&3&(1,1,3,3)&(-4,0,8)\\
\bottomrule
\end{array}
\]
In particular, the displacement depends only on $h(a)$, including for
the two auxiliary letters $D,E$ that share the output letter $3$.
Define the linear map $L:\Q^4\to\Q^3$ by the matrix
\begin{equation}\label{eq:L}
  L=
  \begin{pmatrix}
     2&-2& 0&-4\\
     0& 2&-2& 0\\
     2& 4& 6& 8
  \end{pmatrix}.
\end{equation}
Its columns, in output-letter order $0,1,2,3$, are the four return
displacements above.

\begin{proof}[Proof of Theorem~\ref{thm:main}]
For $\ws=s_0s_1s_2\cdots$, let
\[
  N_0=0,\qquad N_k=\sum_{j=0}^{k-1}|R(s_j)|\quad(k\geq1).
\]
By Lemma~\ref{lem:returns}, these are precisely the zero positions of
$\wt$, in increasing order. The strict increase follows also from the
fact that every return block has positive length.

Write $\ww=w_0w_1w_2\cdots$, where $w_j=h(s_j)$. Additivity of
displacements and the return-block table show that, for all $0\leq p<q$,
\begin{equation}\label{eq:factor-displacement}
  P_{N_q}-P_{N_p}
  =\sum_{j=p}^{q-1}\Delta(R(s_j))
  =L\Par(w_p w_{p+1}\cdots w_{q-1}).
\end{equation}
This holds even though $h$ is not one-to-one, because $R(D)$ and $R(E)$
have equal displacements.

Suppose for a contradiction that a factor of $\ww$ is a weak abelian
cube $xyz$. Let $k_0$ be the position at which it starts and put
\[
  k_1=k_0+|x|,\qquad
  k_2=k_1+|y|,\qquad
  k_3=k_2+|z|.
\]
All three blocks are nonempty, so $k_0<k_1<k_2<k_3$.
There is a vector $f\in\Q_{\geq0}^4$ such that
\[
  \Par(x)=|x|f,\qquad
  \Par(y)=|y|f,\qquad
  \Par(z)=|z|f.
\]
Its coordinates sum to $1$, since $f=\Par(x)/|x|$. If $v=Lf$, then
the third coordinate of $v$ is
\[
  2f_0+4f_1+6f_2+8f_3\geq2>0.
\]
In particular, $v\ne0$. Applying~\eqref{eq:factor-displacement} to
the three blocks gives
\begin{align*}
  P_{N_{k_1}}-P_{N_{k_0}}&=|x|v,\\
  P_{N_{k_2}}-P_{N_{k_1}}&=|y|v,\\
  P_{N_{k_3}}-P_{N_{k_2}}&=|z|v.
\end{align*}
Thus the four distinct vertices $P_{N_{k_0}},P_{N_{k_1}},P_{N_{k_2}},
P_{N_{k_3}}$ lie on the line $P_{N_{k_0}}+\mathbb R v$.
Each associated index is a zero position of $\wt$, contradicting
Proposition~\ref{prop:no-four}. This excludes every weak abelian cube,
regardless of its position or its three block lengths.
\end{proof}

\begin{remark}
There is also a description using only the zero positions of the
four-letter fixed point $\wt$. The return lengths associated with
output letters $0,1,2,3$ are $2,4,6,8$, respectively. Hence, with
$N_0<N_1<\cdots$ enumerating the zero positions of $\wt$,
\[
  w_k=\frac{N_{k+1}-N_k}{2}-1\qquad(k\in\N).
\]
Lemma~\ref{lem:returns} proves that these gaps are always in
$\{2,4,6,8\}$, so this formula is equivalent to the explicit
substitution-and-coding construction in Theorem~\ref{thm:main}.
\end{remark}


\section{Dimension 3}

 Write $e_0=(1,0,0)$, $e_1=(0,1,0)$, and $e_2=(0,0,1)$.
For a word $u=u_0u_1\cdots$ over $\{0,1,2\}$, its walk is
\[
 P_0=0,\qquad P_n=\sum_{j=0}^{n-1}e_{u_j}\quad(n\geq1).
\]
Thus $P_n$ is the Parikh vector of the length-$n$ prefix. Since the sum
of its coordinates is $n$, all walk vertices are distinct.

\subsection{The construction and its precise bound}

Let $\mathcal A=\{A,B,C,D,E\}$. Define the substitution $\tau$ as in Section~\ref{dim4} and
the non-erasing word morphism $g$ by
\[
\begin{array}{c|ccccc}
 a&A&B&C&D&E\\\hline
 \tau(a)&AB&AACA&ADE&AACCE&ADCCA\\
 g(a)&0&1&20&21&21
\end{array}
\]
Both maps extend to finite and infinite words by concatenation. Because
$\tau(A)$ starts with $A$ and every image has length at least two, the
finite words $\tau^k(A)$ are nested prefixes of unbounded length.
They define an infinite fixed point $v=\tau^\infty(A)$.

\begin{theorem}\label{thm:main2}
The infinite ternary word
\[
 u=g\bigl(\tau^\infty(A)\bigr)
\]
defines a walk using only $e_0,e_1,e_2$ with no seven collinear points.
\end{theorem}

The start of the word is
\[
 u=\word{010020001010212101010020001002000100202021021202}\cdots.
\]
 In other words, Theorem
\ref{thm:main2} gives a word with no weak abelian sixth power.

\subsection{A valuation identity for a four-state walk}

Define a three-uniform substitution $\sigma$ on $\{0,1,2,3\}$ by
\[
 \sigma(0)=010,\quad \sigma(1)=232,\quad
 \sigma(2)=101,\quad \sigma(3)=323,
\]
and let $t=\sigma^\infty(0)$. Associate two signs with each state:
\[
\begin{array}{c|rrrr}
 s&0&1&2&3\\\hline
 a(s)&1&1&-1&-1\\
 b(s)&1&-1&1&-1
\end{array}
\]
Write $a_n=a(t_n)$ and $b_n=b(t_n)$. Define
\[
 X_n=\sum_{j<n}a_j,\quad Y_n=\sum_{j<n}b_j,\quad
 R_n=(X_n,Y_n,n).
\]
The substitution gives, for $r\in\{0,1,2\}$,
\begin{equation}\label{eq:sign}
 a_{3n+r}=b_n,\qquad b_{3n+r}=(-1)^r a_n.
\end{equation}
Summing complete three-letter blocks and then a prefix of the next
block gives
\begin{equation}\label{eq:prefix}
 X_{3n+r}=3Y_n+r b_n,\qquad
 Y_{3n+r}=X_n+\epsilon_r a_n,
 \qquad(\epsilon_0,\epsilon_1,\epsilon_2)=(0,1,0).
\end{equation}

For a nonzero integer $z$, let $\vp(z)$ be the largest integer $j\geq0$
such that $3^j$ divides $z$. The sign of $z$ is irrelevant to this
definition. Set $Q(X,Y)=X^2-3Y^2$.

\begin{lemma}\label{lem:valuation2}
If $0\leq m<n$ and $t_m=t_n$, then $Q(X_n-X_m,Y_n-Y_m)$ is nonzero,
and
\[
 \vp\bigl(Q(X_n-X_m,Y_n-Y_m)\bigr)=\vp(n-m).
\]
\end{lemma}
\begin{proof}
Write $m=3p+r$ and $n=3q+s$, where $r,s\in\{0,1,2\}$.
If $r\ne s$, equality of the endpoint states implies equality of their
first signs. Equation~\eqref{eq:sign} therefore gives $b_p=b_q$.
By~\eqref{eq:prefix},
\[
 X_n-X_m\equiv (s-r)b_p\not\equiv0\pmod3.
\]
Consequently $Q(X_n-X_m,Y_n-Y_m)\not\equiv0\pmod3$, and both
valuations in the claimed identity are zero.

If $r=s$, the two signs in~\eqref{eq:sign} show that $a_p=a_q$ and
$b_p=b_q$, hence $t_p=t_q$. The offsets in~\eqref{eq:prefix} cancel,
so
\[
 X_n-X_m=3(Y_q-Y_p),\qquad Y_n-Y_m=X_q-X_p.
\]
It follows that
\[
 Q(X_n-X_m,Y_n-Y_m)
 =-3Q(X_q-X_p,Y_q-Y_p),\qquad n-m=3(q-p).
\]
Thus the desired identity descends to the equal-state pair $(p,q)$,
adding one to each valuation. After exactly $\vp(n-m)$ descents, the
two indices have different residues modulo three, so the preceding
base case applies. The same argument proves that the quadratic-form
value is nonzero throughout.
\end{proof}

\begin{lemma}\label{lem:four}
For each fixed state $s\in\{0,1,2,3\}$, the set
$\{R_n:t_n=s\}$ contains no four collinear vertices.
\end{lemma}
\begin{proof}
Suppose four such vertices are collinear. As they are distinct integer
points, there are a primitive integer direction
$D=(D_X,D_Y,D_H)$ with $D_H>0$, an integer point $R$, and distinct
integers $q_0,q_1,q_2,q_3$ such that the four vertices are $R+q_iD$.
To justify the integrality of the parameters, choose $D$ by dividing
one nonzero displacement by the gcd of its three coordinates. Because
the coordinates of $D$ have gcd one, an integer linear combination
of them is one; applying that combination to any other integer
displacement parallel to $D$ gives its integer parameter.

Lemma~\ref{lem:valuation2} implies $Q(D_X,D_Y)\ne0$. Applying that
lemma to any pair of these four vertices and using homogeneity gives
\[
 2\vp(q_j-q_i)+\vp\bigl(Q(D_X,D_Y)\bigr)
 =\vp(q_j-q_i)+\vp(D_H).
\]
Therefore all six pairwise differences $q_j-q_i$ have the same
valuation, say $h$. Every $q_i-q_0$ is divisible by $3^h$, and the
four integers $(q_i-q_0)/3^h$ must be pairwise distinct modulo three.
There are only three residue classes, a contradiction.
\end{proof}

\subsection{Return words and a walk with diagonal steps}

Consider the five finite words
\[
\begin{array}{c|ccccc}
 a&A&B&C&D&E\\\hline
 r(a)&01&0232&013231&02323231&01323232
\end{array}
\]
Each starts with a single zero and contains no further zero. Direct
concatenation verifies the five exact identities
\begin{align*}
 \sigma(r(A))&=r(A)r(B),\\
 \sigma(r(B))&=r(A)r(A)r(C)r(A),\\
 \sigma(r(C))&=r(A)r(D)r(E),\\
 \sigma(r(D))&=r(A)r(A)r(C)r(C)r(E),\\
 \sigma(r(E))&=r(A)r(D)r(C)r(C)r(A).
\end{align*}
Equivalently, $\sigma\circ r=r\circ\tau$. Hence $r(v)$ is a fixed
point of $\sigma$ starting with zero, and it must equal $t$: any such
fixed point has prefix $\sigma^k(0)$ for every $k$, and these prefixes
have lengths tending to infinity. Thus the boundaries of the return
words $r(v_0),r(v_1),\ldots$ are exactly the positions at which the
state of $t$ is zero.

Let
\[
 \ell_0=0,\qquad \ell_n=\sum_{j<n}|r(v_j)|.
\]
Then $t_{\ell_n}=0$. The displacements of $R$ across the five return
words, divided by two, are
\[
\begin{array}{c|ccccc}
 a&A&B&C&D&E\\\hline
 \tfrac12\Delta R(r(a))
 &(1,0,1)&(-1,1,2)&(0,-1,3)&(-2,0,4)&(-2,0,4).
\end{array}
\]
These entries follow by adding the sign pairs of the displayed words
and taking half their lengths for the third coordinate.

Now define a walk $S_0,S_1,\ldots$ with $S_0=0$ and increments
\[
 d(A)=e_0,\quad d(B)=e_1,\quad
 d(C)=e_0+e_2,\quad d(D)=d(E)=e_1+e_2,
 \qquad S_{n+1}=S_n+d(v_n).
\]
The integer matrix
\[
 M=\begin{pmatrix}1&-1&-1\\0&1&-1\\1&2&2\end{pmatrix},
 \qquad \det M=6,
\]
maps each of these four increment types to its entry in the return
displacement table. It follows by summing increments that
\[
 MS_n=\tfrac12 R_{\ell_n}\quad\hbox{for every }n\geq0.
\]
Since $M$ is invertible, collinearity of $S$ vertices is equivalent to
collinearity of the corresponding return vertices. Lemma~\ref{lem:four}
therefore proves that every affine line contains at most three points
of the vertex set
\[
 \mathcal S=\{S_n:n\geq0\}.
\]

\subsection{Subdivision gives the seven-point bound}

The coding $g$ subdivides every diagonal increment of $S$ into its two
positive coordinate steps, taking the $e_2$ step first. In detail,
$C$ contributes $e_2,e_0$, while $D$ and $E$ contribute $e_2,e_1$.
An $A$ or $B$ increment is left as a single unit step.

Consequently every vertex of the unit-step walk $P$ is either an
original vertex $S_n$ or an inserted vertex $S_n+e_2$. In particular,
\[
 \{P_k:k\geq0\}\subseteq \mathcal S\cup(\mathcal S+e_2).
\]
Translation preserves collinearity, so every line contains at most
three vertices of $\mathcal S+e_2$, as well as at most three vertices
of $\mathcal S$. Thus it contains at most six vertices of $P$.
This proves the avoidance assertion in Theorem~\ref{thm:main2}.

\section{Declaration of AI usage}

\begin{thebibliography}{99}

\bibitem{AP16}
Sergey Avgustinovich and Svetlana Puzynina.
\newblock Weak abelian periodicity of infinite words.
\newblock {\it Theory of Computing Systems} \textbf{59} (2016), 161--179.

\bibitem{Bro71a}
T. C. Brown.
\newblock Is there a sequence on four symbols in which no two adjacent
segments are permutations of one another?
\newblock {\it Amer. Math. Monthly} {\bf 78} (1971), 886--888.

\bibitem{Bro71b}
T. C. Brown.
\newblock Advanced problem 5811.
\newblock {\it Amer. Math. Monthly} {\bf 78} (1971), 798.

\bibitem{CK}
Stijn Cambie and Erik Kalviainen.
\newblock An infinite small-step $\Z^3$-walk with no collinear triple.
\newblock Arxiv preprint arXiv:2609.01766 [math.CO],
September 1 2026.  Available at
\url{https://arxiv.org/abs/2609.01766v1}.

\bibitem{FP23}
Gabriele Fici and Svetlana Puzynina.
\newblock Abelian combinatorics on words: a survey.
\newblock {\it Computer Science Review} \textbf{47} (2023), article~100532.

\bibitem{GR79}
Joseph L.~Gerver and L.~Thomas Ramsey.
\newblock On certain sequences of lattice points.
\newblock {\it Pacific Journal of Mathematics} \textbf{83} (1979), 357--363.

\bibitem{Ko26a}
Samuel Korsky.
\newblock North-east lattice paths with few collinear vertices.
\newblock Arxiv preprint arXiv:2607.02832 [math.CO], July 2 2026.
\newblock Available at \url{https://arxiv.org/abs/2607.02832}.

\bibitem{Ko26b}
Samuel Korsky.
\newblock Long lattice paths with no three collinear vertices.
\newblock Arxiv preprint arXiv:2608.07906 [math.CO], August 8 2026.
\newblock Available at \url{https://arxiv.org/abs/2608.07906}.

\bibitem{Lid24}
Thomas F.~Lidbetter.
\newblock Improved bound for the Gerver--Ramsey collinearity problem.
\newblock {\it Discrete Mathematics} \textbf{347} (2024), article~113718.

\bibitem{Mo72}
P. L. Montgomery. 
\newblock Solution to Advanced Problem 5811: 
collinear points on a monotonic polygon.
\newblock {\it American Mathematical Monthly} \textbf{79} (1972), 1143--1144.

\bibitem{OEIS}
N. J. A. Sloane et al.
\newblock The On-Line Encyclopedia of Integer Sequences.
\newblock Electronic resource, 2026.  Available at
\url{https://oeis.org}.





\end{thebibliography}

\end{document}